\chapter{Experimental Design}
\label{ch:experiments}

\section{Overview}

All experiments evaluate the six algorithmic strategies: State-Expanded Dijkstra (B1),
Greedy (B2), DP (B3), Standard A* (B4), RF-A* (B5), and PF-A* (proposed), on the
variable-purchase model.
Every algorithm receives the same inputs: graph $G$, prices $p$, initial fuel $f_0$, tank
capacity $C$, and efficiency $\eta$.
The primary research question is whether PF-A* achieves the same optimal cost as
Dijkstra/DP/A*/RF-A* while expanding fewer states, and whether all optimal algorithms
outperform Greedy in cost.

All experiments are conducted on Microsoft Azure Standard\_D4s\_v5 virtual machines
(4 vCPUs Intel Xeon Platinum 8370C 2.8\,GHz, 16\,GB DDR4, Ubuntu 22.04\,LTS,
Python~3.11, NetworkX~3.x, OSMnx).
Runtime is measured with \texttt{time.perf\_counter\_ns()} at the algorithmic level,
excluding I/O.
Statistical significance uses the Wilcoxon signed-rank test ($\alpha=0.05$, Bonferroni
correction for multiple comparisons).
Two run budgets are used in the final report.
The synthetic experiments (Experiments~1--5 and~7) use a reduced but reproducible
benchmark profile of 3 timed iterations per configuration, summarised in
Section~\ref{sec:current_synth_snapshot}; the full target profile of 10{,}000 timed runs
per configuration is reported as future work.
The real-route evaluation (Experiment~6) uses the full 10{,}000 timed runs per
algorithm--corridor combination, producing 210{,}000 result rows in total.

\section{Baselines}

\textbf{State-Expanded Dijkstra (B1)} is the cost lower bound: it exhaustively explores
$G^{*}$ and is guaranteed to find the minimum-cost variable-purchase strategy.
Dynamic Programming (B3) and PF-A* are treated as the exact-agreement set with
Dijkstra: all three are expected to return identical costs on every instance, and any
divergence among these three flags a correctness bug.
Standard A* (B4) and RF-A* (B5) inherit admissibility from their heuristics in theory but,
in the current implementation, exhibit small positive cost gaps against Dijkstra on some
instances (see Chapter~\ref{ch:results}).
These gaps are tracked as an implementation-audit item rather than treated as exact
agreement.

\textbf{Full-tank-only Greedy (B2)} is the cost upper bound: it fills to capacity at every
station regardless of price, representing worst-case uninformed driver behavior.

\textbf{RF-A* (B5)} is the direct algorithmic baseline for PF-A*: both are optimal heuristic
search algorithms with the same worst-case complexity, operating on the same state-expanded
graph under the variable-purchase model.
The sole difference is the heuristic. $h_{\mathrm{RF}}(v)$ is fuel-level agnostic;
$h_{\mathrm{PF}}(v,f)$ is fuel-level aware.
Experiment~5 isolates this difference directly.

\section{Evaluation Metrics}

\begin{table}[ht]
\centering
\caption{Evaluation metrics.}
\label{tab:metrics}
\begin{tabular}{lll}
\toprule
Metric & Unit & Purpose \\
\midrule
Total fuel cost      & THB   & Primary objective; must match B1 for optimal algorithms \\
Runtime              & ms    & Computational efficiency \\
Nodes expanded       & count & Search space explored; key metric for PF-A* vs.\ RF-A* \\
Memory consumption   & MB    & Scalability of $O(VF)$ state space \\
Refuelling stops     & count & Practical convenience; tracks partial-fill decisions \\
Sub-optimality gap    & \% vs.\ B1 & Cost excess of Greedy (correctness check for others) \\
\bottomrule
\end{tabular}
\end{table}

\section{Experimental Configurations}

\subsection{Experiment 1: Scalability (Graph Size)}

Fix $\sigma=10\%$, $C=40$\,L, $f_0=0.7C=28$\,L; vary graph size.

\begin{center}
\begin{tabular}{llll}
\toprule
Size & $|V|$ & $|E|$ (approx.) & Runs \\
\midrule
Small  & 100    & 500    & 3 \\
Medium & 1,000  & 5,000  & 3 \\
Large  & 10,000 & 50,000 & 3 \\
\bottomrule
\end{tabular}
\end{center}

\textit{Hypothesis:} Greedy scales as $O(V^{2})$; all other algorithms scale as
$O((EF+VF^{2})\log(VF))$.
PF-A* expands fewer nodes than RF-A* at all sizes, with the gap growing as density
increases.

\subsection{Experiment 2: Fuel Price Variance}

Fix $|V|=1{,}000$, $C=40$\,L, $f_0=28$\,L; vary $\sigma \in \{0\%,5\%,10\%,15\%,20\%\}$.

\textit{Hypothesis:} At $\sigma=0\%$ (uniform prices), all algorithms produce identical costs
and PF-A* offers no node-expansion advantage over RF-A*.
As $\sigma$ increases, optimal algorithms save 5--15\% more cost versus Greedy, and PF-A*
expands progressively fewer states than RF-A* because
$p_{\min}^{\mathrm{reach}}(v,f) \gg p_{\min}$ more often.

\subsection{Experiment 3: Tank Capacity and Initial Fuel}

Fix $|V|=1{,}000$, $\sigma=10\%$; vary $C \in \{20,40,60,80\}$\,L with $f_0=0.7C$.

\textit{Hypothesis:} Smaller $C$ tightens the fuel constraint, increasing the importance of
optimal purchase decisions and widening the cost gap between Greedy and optimal algorithms.
PF-A*'s advantage over RF-A* is largest at small $C$.

\subsection{Experiment 4: Variable Initial Fuel Level}

Fix $|V|=1{,}000$, $C=40$\,L, $\sigma=10\%$; vary
$f_0 \in \{0.25C, 0.50C, 0.75C, 1.0C\}$.

\textit{Hypothesis:} Lower $f_0$ forces earlier mandatory refueling decisions.
PF-A*'s heuristic is most informative at low $f_0$ (where $\max(d_{\mathrm{geo}}/\eta - f, 0)$
is large), so its node-expansion advantage over RF-A* should grow as $f_0$ decreases.

\subsection{Experiment 5: RF-A* vs.\ PF-A* Node Expansions (Primary Novelty Experiment)}

Fix $|V|=1{,}000$; jointly vary $\sigma \in \{5\%,10\%,20\%\}$ and station density
$\in \{0.5, 1, 2, 5\}$ per 100\,km.
Run only RF-A* and PF-A*.
Primary metric: nodes expanded.
Secondary check: both must return identical optimal costs on every instance.

\textit{Hypothesis:} PF-A* expands fewer nodes than RF-A* in all configurations with
$\sigma > 0\%$.
The reduction is largest when station density is high (many stations reachable per state, so
$p_{\min}^{\mathrm{reach}}$ differs most from global $p_{\min}$) and price variance is high.

\subsection{Experiment 6: Real Thai Routes}

Experiment~6 is designed around the actual curated real-route dataset already available in the
repository, rather than around a future nationwide OSM benchmark.
The dataset contains three fixed Thai corridor instances:
Bangkok $\to$ Chiang Mai, Bangkok $\to$ Phuket, and Bangkok $\to$ Khon Kaen.
Each instance is a directed corridor graph with explicit source and destination anchors,
intermediate refueling nodes, precomputed feasible travel legs under a 400\,km full-tank
range, and province-level fuel prices tied to the 3~April~2026 Gasohol~91 snapshot
documented in \texttt{data/real\_routes/corridors.json}.

All three corridors use the same vehicle parameters as the synthetic study:
$C=40$\,L, $f_0=28$\,L, and $\eta=10$\,km/L.
The current generated instances contain 27 stations and 522 directed edges for Chiang Mai,
31 stations and 624 directed edges for Phuket, and 18 stations and 286 directed edges for
Khon Kaen.
Each route is evaluated under the variable-purchase model, with a full-tank-only Dijkstra run
used as the practical baseline.

\textit{Purpose:} Validates whether the synthetic findings transfer to realistic Thai corridor
instances.
The primary empirical questions are whether PF-A* still matches the exact optimum while
expanding fewer states, and whether variable purchase remains cheaper than full-tank-only
routing on all three real-route corridors.

\textit{Current design note:} Because the available instances are curated corridor abstractions
rather than raw OSM forecourt graphs, Experiment~6 should be written as a real-route corridor
case study.
That framing is strong enough for the available data and avoids overstating station-level
realism.

\subsection{Experiment 7: Variable Purchase vs.\ Full-Tank-Only Model Comparison}

Fix $|V|=1{,}000$, $f_0=0.7C$; jointly vary $\sigma \in \{0\%,5\%,10\%,20\%\}$ and
$C \in \{20,40,60,80\}$\,L.
Run both models: full-tank-only Dijkstra on $G$ and variable-purchase Dijkstra (B1) on
$G^{*}$.
Primary metric: cost gap (THB and \%).
Secondary metric: runtime overhead of the variable-purchase model.

\textit{Hypothesis 1 (cost):} Variable purchase strictly dominates full-tank-only in cost on
all instances with $\sigma > 0\%$.
The gap grows with $\sigma$ and shrinks with $C$.

\textit{Hypothesis 2 (runtime overhead):} Variable-purchase Dijkstra is slower than
full-tank-only Dijkstra by a factor of $O(F\log F)$ in theory, expected to be 10--40$\times$
in practice for $C=40$\,L, $\varepsilon=1$\,L.
PF-A* is expected to close this gap significantly.

\textit{Hypothesis 3 (tank size interaction):} At $C=20$\,L the variable-purchase model
saves the most; at $C=80$\,L the savings diminish.
This interaction will be visualised as a $\sigma \times C$ heatmap of cost savings.

\section{Correctness Verification}

Before benchmarking, correctness is validated independently:
\begin{itemize}
    \item \textbf{Unit tests:} each algorithm tested on 50 hand-crafted small graphs with
          known optimal solutions, including the running example
          ($f_0=28$\,L, optimal $=216$\,THB).
    \item \textbf{Cross-validation (exact-agreement set):} on all synthetic instances,
          Dijkstra (B1), DP (B3), and PF-A* must agree on optimal cost within
          floating-point tolerance.
    \item \textbf{A*/RF-A* audit:} Standard A* and RF-A* are cross-checked against
          Dijkstra; any residual cost gap is logged as an implementation-audit item
          rather than suppressed.
    \item \textbf{Feasibility checks:} every returned solution is verified to never drop below
          0\,L fuel at any point along the route.
    \item \textbf{Infeasibility detection:} instances with no valid path are tested to confirm
          all algorithms correctly return ``infeasible.''
\end{itemize}
